% Part: history
% Chapter: biographies 
% Section: georg-cantor

\documentclass[../../../include/open-logic-section]{subfiles}

\begin{document}

\olfileid[pt]{his}{bio}{can}

\olsection{Georg Cantor}

\olphoto{cantor-georg}{Georg Cantor}

Uma biografia inicial de Georg Cantor (\textsc{gay}-org
\textsc{kahn}-tor) afirmava que ele tinha nascido e sido encontrado
num navio que seguia para São Petersburgo, na Rússia, e que seus
pais eram desconhecidos. Isto, porém, não é verdade; embora ele
tivesse nascido em São Petersburgo em 1845.

Cantor se doutorou em matemática na Universidade de Berlim em
1867. É conhecido por seu trabalho em teoria dos conjuntos e é
considerado o fundador da teoria dos conjuntos como disciplina de
pesquisa distinta.
Foi o primeiro a provar que existem conjuntos infinitos de tamanhos
diferentes. Suas teorias, e sobretudo a teoria dos infinitos,
provocaram muito debate entre os matemáticos da época, e seu
trabalho foi controverso.

As convicções religiosas de Cantor e seu trabalho matemático
estavam inextricavelmente ligados; ele chegou a afirmar que a teoria
dos números transfinitos lhe tinha sido comunicada diretamente por
Deus. Mais tarde, Cantor sofreu de doença mental. A partir de 1894, e
com maior frequência nos últimos anos, Cantor foi internado. As fortes
críticas a seu trabalho, incluindo um rompimento com o matemático
Leopold Kronecker, levaram à depressão e à perda de interesse pela
matemática. Durante episódios depressivos, Cantor se dedicava à
filosofia e à literatura, e chegou a publicar uma teoria de que
Francis Bacon era o autor das peças de Shakespeare.

Cantor morreu em 6 de janeiro de 1918, num sanatório em Halle.

\begin{reading} 
Para biografias completas de Cantor, ver \citet{Dauben1990} e
\citet{Grattan-Guinness1971}.  As visões radicais de Cantor são
também descritas no programa da BBC Radio 4 \emph{A Brief History of
  Mathematics} \citep{Sautoy2014}.  Se quiser ouvir as teorias de
Cantor em forma de rap, ver \citet{Rose2012}.
\end{reading}

\end{document}
